% Data science ideas: written by hand, read by docs/make_manual.py (% Data science ideas: written by hand, read by docs/make_manual.py (% Data science ideas: written by hand, read by docs/make_manual.py (% Data science ideas: written by hand, read by docs/make_manual.py (\input{ideas}).
% Plain words first, then the formal name, so a learner can follow a paper's Methods section.

\unchapter{Describing numbers}{ideas:describe}

Before any test or model, a few numbers describe a column well: where its middle is, and how
spread out it is. simplydata's summary words say each one in a sentence.

\unsection{Average or median?}{ideas:average-median}

The \textbf{average} (formally the \textbf{arithmetic mean}, written $\bar{x}$) adds all the values
and divides by how many there are:
\[
  \bar{x} = \frac{x_1 + x_2 + \dots + x_n}{n}.
\]
The \textbf{median} is the middle value when the values are put in order (with an even number of
values, the average of the two in the middle). Half the values are below it and half above.

Both describe the ``middle'', but they react differently to values far from the rest
(\textbf{outliers}). Five children earn pocket money of 2, 3, 3, 4 and 50 a week: the average is
12.4, though four of the five get 4 or less; the median is 3. One extreme value pulls the average
towards it but hardly moves the median. So:

\begin{itemize}
  \item Use the \textbf{average} when the values are roughly symmetric (a bell shape), and when you
    want to add things up (total = average $\times$ count).
  \item Use the \textbf{median} when the values are lopsided (\textbf{skewed}), such as incomes,
    house prices, waiting times or bacterial counts, or when there are outliers you cannot
    remove.
  \item When the two are far apart, say so: it tells the reader the values are skewed.
\end{itemize}

In simplydata: \code{table.average\_of("marks")}, \code{table.median\_of("marks")},
\code{nums.average()}, \code{nums.median()}, and on charts \code{show="average"} or
\code{show="median"}. Lesson~\ref{lesson:5} draws both.

\unsection{Spread: range, percentiles and standard deviation}{ideas:spread}

Two classes can have the same average mark, one with every child near 65 and the other with
marks from 20 to 100. The \textbf{spread} tells them apart.

\begin{itemize}
  \item The \textbf{range} is the biggest value minus the smallest. It is easy to understand but
    depends on just two values, often the strangest ones.
  \item A \textbf{percentile} is the value that a given share of the values is below:
    \code{nums.percentile(90)} is the value that 90\,\% are below. The median is the 50th
    percentile. The middle half of the values lies between the 25th and 75th percentiles; that
    distance is the \textbf{interquartile range (IQR)}, and it is the box in a box plot
    (\code{chart.boxes}).
  \item The \textbf{standard deviation} (SD, written $s$) is about how far a typical value is from
    the average:
    \[
      s = \sqrt{\frac{\sum_{i=1}^{n} (x_i - \bar{x})^2}{n - 1}}.
    \]
    Its square, $s^2$, is the \textbf{variance}. Dividing by $n-1$ instead of $n$ (Bessel's
    correction) makes the sample variance an unbiased estimate of the variance of everything the
    sample came from; simplydata always divides by $n-1$, as pandas and R do (NumPy divides by $n$
    unless told otherwise).
\end{itemize}

The SD goes with the average and the IQR with the median. For all the marks of the practice table:
``average 64.9 (SD 15.0)'' or ``median 68.5 (IQR 54.3 to 74.8)''.

\unsection{The shape of the values: distributions}{ideas:distribution}

The way values are spread out is their \textbf{distribution}; a histogram
(\code{chart.histogram}) or a density curve (\code{chart.density}) shows it. Many measurements,
such as heights or measuring errors, follow a bell shape called the \textbf{normal} (or
Gaussian) \textbf{distribution}. For a normal distribution about 68\,\% of the values lie within
one SD of the average and about 95\,\% within two SD.

Not everything is bell-shaped. Counts and concentrations are often skewed to the right (many
small values, a few large ones); a log scale (\code{chart.log\_scale\_y()}) often makes them look
more symmetric. Look at the distribution before choosing what to report and which test to use.

\ideachapter{How sure are we?}{ideas:uncertainty}

We almost never measure everything. We measure a \textbf{sample} (18 children, 40 water samples)
and want to learn about the \textbf{population} it came from (all children of that age, the whole
fjord). A number worked out from a sample, such as its average, is an \textbf{estimate}, and
another sample would give a slightly different one. These ideas say how much.

\unsection{Standard deviation or standard error?}{ideas:sd-se}

They sound alike and are often mixed up, but they answer different questions.

\begin{center}\small
\begin{tabularx}{\textwidth}{@{}>{\raggedright\arraybackslash}p{0.19\textwidth}>{\raggedright\arraybackslash}X>{\raggedright\arraybackslash}X@{}}
\toprule
 & \textbf{Standard deviation (SD)} & \textbf{Standard error of the mean (SE, SEM)}\\
\midrule
Describes & the \emph{values}: how spread out they are & the \emph{average}: how precisely the sample pins it down\\
Formula & $s$ (above) & $\mathrm{SE} = s / \sqrt{n}$\\
More data & stays about the same & gets smaller (4 times the data halves it)\\
Use it to & show the spread of individuals & compare averages, build confidence intervals\\
simplydata & \code{error="standard deviation"} & \code{error="standard error"}\\
\bottomrule
\end{tabularx}
\end{center}

Error bars of the SE are always shorter than those of the SD, and can look reassuringly small. That
is why a figure caption must say which one was drawn; simplydata prints a line saying it, ready
for your caption.

\unsection{Confidence intervals}{ideas:ci}

A \textbf{95\,\% confidence interval} (CI) is a range of plausible values for the true number in
the population, worked out from the sample. For an average it is
\[
  \bar{x} \pm t_{0.975,\,n-1} \times \mathrm{SE},
\]
where $t_{0.975,\,n-1}$ is a number from the $t$ distribution: about 2 for most samples, a bit
more for small ones (2.26 for $n = 10$, 2.09 for $n = 20$).

What ``95\,\%'' means, exactly: if you repeated the whole study many times and built an interval
each time, about 95 out of 100 of those intervals would contain the true value. It does
\emph{not} mean that the true value moves around, nor that 95\,\% of the individual values lie
inside (that would be about $\bar{x} \pm 2\,\mathrm{SD}$ for a bell shape).

\begin{itemize}
  \item A \textbf{narrow} interval means the data pin the value down well; a \textbf{wide} one
    means you need more data.
  \item An interval says more than a p-value: it gives the size of the effect \emph{and} how sure
    you are of it. \code{table.fit\_straight\_line} prints the slope with its 95\,\% CI.
  \item On charts: \code{error="confidence interval"} on bars and lines, and
    \code{fit\_line(confidence\_band=True)} for a fitted line.
\end{itemize}

\unsection{Why averages behave well: the central limit theorem}{ideas:clt}

Even when the values themselves are skewed, the \emph{average} of many of them is close to
bell-shaped; this is the \textbf{central limit theorem}. It is why tests on averages (t-tests,
ANOVA) usually work well with 30 or more values per group, even for values that are not
bell-shaped. With small groups of skewed values, rank-based tests are safer
(\code{based\_on="ranks"}).

\ideachapter{Testing: real difference, or chance?}{ideas:testing}

\unsection{The null hypothesis and the p-value}{ideas:p-value}

Two groups almost never have exactly the same average, even when nothing real separates them:
samples differ by chance. A statistical test asks whether a difference is bigger than chance alone
would often produce. It works like a court case, where the accused is innocent until proven
guilty:

\begin{enumerate}
  \item Assume there is \textbf{no real difference}. This assumption is the \textbf{null
    hypothesis} ($H_0$). The \textbf{alternative hypothesis} ($H_1$) is that there is one.
  \item Work out how surprising your data would be if $H_0$ were true. The \textbf{p-value} is the
    probability of a difference at least as big as the one you saw, if there were really no
    difference.
  \item A small p-value means the data are hard to explain by chance alone. By custom, below
    $0.05$ is called \textbf{statistically significant}; that threshold is the
    \textbf{significance level} $\alpha$.
\end{enumerate}

simplydata says the p-value in words: below 0.001 ``very unlikely to be chance'', below 0.05
``unlikely to be chance'', otherwise ``could easily be chance''.

\begin{goodtoknow}
What a p-value is \textbf{not}:
\begin{itemize}
  \item It is \textbf{not} the probability that $H_0$ is true, nor the probability that the result
    is a fluke.
  \item A small p-value does \textbf{not} say the difference is big or important. With thousands
    of values, tiny, meaningless differences become ``significant''. The \emph{effect size} says
    how big it is.
  \item A large p-value does \textbf{not} show that there is no difference. It only says these
    data cannot tell. ``Absence of evidence is not evidence of absence.''
  \item 0.049 and 0.051 are practically the same evidence. Report the exact p-value, not just
    ``$p < 0.05$''.
\end{itemize}
\end{goodtoknow}

\unsection{Two kinds of mistake, and power}{ideas:errors}

\begin{itemize}
  \item A \textbf{false alarm} (\textbf{type I error}, a false positive): the test says there is a
    difference when there is none. With $\alpha = 0.05$ this happens in about 1 test in 20 when
    $H_0$ is true.
  \item A \textbf{missed difference} (\textbf{type II error}, a false negative): there is a real
    difference, but the test does not find it. Small samples miss differences often.
\end{itemize}

The \textbf{power} of a study is the chance of finding a difference of a given size if it is
really there. More data, less noisy measurements and bigger true differences all raise power. A
study with low power can only find large effects; plan the sample size before collecting data.

\unsection{Many tests at once}{ideas:multiple}

Every test at $\alpha = 0.05$ carries a 5\,\% risk of a false alarm, and the risks add up: with 20
tests on pure noise, you expect about one ``significant'' result. Correcting for
\textbf{multiple comparisons} keeps the chance of \emph{any} false alarm at 5\,\%.
\code{table.compare\_each\_pair} uses \textbf{Holm's method}, which is never less powerful than the
simpler Bonferroni correction (multiply each p-value by the number of tests). Decide which
comparisons you will make before you look at the data.

\unsection{Effect sizes: how big is the difference?}{ideas:effect-size}

The p-value answers ``could it be chance?''. The \textbf{effect size} answers ``how big is it?'',
which is usually what matters. Report both, ideally with a confidence interval for the effect.

\begin{center}\small
\begin{tabularx}{\textwidth}{@{}>{\raggedright\arraybackslash}p{0.26\textwidth}>{\raggedright\arraybackslash}X>{\raggedright\arraybackslash}p{0.25\textwidth}@{}}
\toprule
\textbf{Effect size} & \textbf{What it says} & \textbf{Given by}\\
\midrule
Difference in averages & in the units of the data (``7.6 marks higher''): often the clearest & \code{compare\_averages}, \code{compare\_each\_pair}\\
Cohen's $d$, Hedges' $g$ & the difference in standard deviations; $g$ corrects $d$ for small samples. Rough guide: 0.2 small, 0.5 medium, 0.8 large & \code{compare\_averages} (2 groups), \code{compare\_average\_to}, \code{compare\_before\_after} ($d_z$)\\
Omega-squared $\omega^2$ & the share of all the differences that the groups explain (0 to 1) & \code{compare\_averages} (3+ groups)\\
Cramér's $V$ & how strongly two categories are linked, 0 (none) to 1 & \code{compare\_counts}\\
Correlation $r$ & how closely two columns go up and down together, $-1$ to $1$ & \code{correlation\_between}\\
$R^2$ & the share of the differences in $y$ that a line or curve explains & \code{fit\_straight\_line}, \code{fit\_formula}\\
\bottomrule
\end{tabularx}
\end{center}

The ``small, medium, large'' labels (from Jacob Cohen) are only a last resort: a ``small''
difference in survival can be important, and a ``large'' one in a lab assay may not be.

\ideachapter{Links between columns}{ideas:links}

\unsection{Correlation}{ideas:correlation}

A \textbf{correlation} measures how closely two columns of numbers go up and down together, from
$-1$ (one always goes down when the other goes up) through $0$ (no straight-line link) to $1$.

\begin{itemize}
  \item \textbf{Pearson's $r$} (the usual one) measures a \emph{straight-line} link. It is
    sensitive to outliers, and can be near 0 for a strong curved link (a U shape).
  \item \textbf{Spearman's $\rho$} (rho) uses ranks: it measures any link that keeps going one way
    (always up, or always down), and is robust to outliers.
    \code{table.correlation\_between("hours", "marks", method="spearman")}.
  \item $r^2$ is the share of the differences in one column that a straight line through the other
    explains. $r = 0.5$ explains only a quarter.
\end{itemize}

Always draw the dots (\code{chart.dots}) before trusting a correlation: very different patterns
can have the same $r$ (Anscombe's quartet is the famous example).

\unsection{Correlation is not cause}{ideas:cause}

Ice cream sales and drownings go up together, but ice cream does not cause drowning: hot weather
drives both. When two things go together there are several possible reasons:

\begin{enumerate}
  \item A causes B.
  \item B causes A (\textbf{reverse causation}): do people exercise because they are healthy, or
    are they healthy because they exercise?
  \item A third thing causes both (a \textbf{confounder}): hot weather, age, wealth, or the size of
    a hospital.
  \item Chance: test enough pairs of columns and some will line up by luck.
\end{enumerate}

Only an \textbf{experiment} where you change A yourself, ideally by random assignment (a
randomised controlled trial), shows cause directly. With observational data, say ``is associated
with'', not ``causes''. simplydata's link words remind you: ``A link does not show that one thing
causes the other.''

\unsection{Fitting a line or a formula}{ideas:regression}

\textbf{Linear regression} finds the straight line $y = a + b\,x$ that goes as close as possible to
the dots, by making the sum of the squared vertical distances from the dots to the line as small as
possible (\textbf{least squares}). The \textbf{slope} $b$ says how much $y$ changes for each 1 more
of $x$; its p-value tests whether the slope could be 0, and its confidence interval gives a range of
plausible slopes. $R^2$ says how much of the differences in $y$ the line explains.
\code{table.fit\_straight\_line(x="hours", y="marks")} prints all of these;
\code{chart.fit\_line()} draws the line.

When the relation is curved, such as growth that multiplies by the same share each year, write the
formula yourself: \code{table.fit\_formula("population = a * exp(b * (year - 2000))")} finds $a$ and
$b$ by \textbf{nonlinear least squares}. A fitted line or curve describes the data you have;
predicting far outside their range (\textbf{extrapolation}) is risky.

\ideachapter{The shape of a table}{ideas:tables}

\unsection{Tidy data: long and wide tables}{ideas:tidy}

The same marks can be stored in two shapes:

\begin{center}\small
\begin{tabular}{@{}lll@{\hspace{3em}}llll@{}}
\multicolumn{3}{@{}l}{\textbf{Long (tidy)}} & \multicolumn{4}{l@{}}{\textbf{Wide}}\\
\toprule
name & subject & marks & name & Maths & Science & English\\
\midrule
Asha & Maths & 61 & Asha & 61 & 66 & 77\\
Asha & Science & 66 & Ben & 73 & 98 & 71\\
Asha & English & 77 & & & & \\
Ben & Maths & 73 & & & & \\
\multicolumn{3}{@{}l}{\dots} & & & & \\
\bottomrule
\end{tabular}
\end{center}

A table is \textbf{tidy} (Hadley Wickham's term) when each \textbf{variable} is a column, each
\textbf{observation} is a row, and each value has its own cell. The long table is tidy: it has one
column for what was measured (\code{marks}) and columns that say which child and which subject.

\begin{itemize}
  \item \textbf{Long tables} suit charts, groups and tests: \code{chart.bars(x="subject", y="marks")}
    and \code{compare\_averages("marks", between="subject")} need the subject as a column.
  \item \textbf{Wide tables} suit sums across columns (a total per child) and paired tests
    (\code{compare\_before\_after("Maths", "Science")}), and are how people often type data into
    a spreadsheet.
  \item \code{table.turn\_columns\_into\_rows} makes a wide table long (in pandas: \code{melt}; in R:
    \code{pivot\_longer}); \code{table.turn\_rows\_into\_columns} makes a long table wide (\code{pivot}).
\end{itemize}

\unsection{Missing values and imputation}{ideas:missing}

Real data has empty cells: a sample was lost, a question skipped. What to do depends on
\emph{why} they are empty:

\begin{itemize}
  \item \textbf{Missing completely at random}: the reason has nothing to do with the data (a tube
    dropped). Leaving those rows out loses some power but does not bias the result.
  \item \textbf{Missing at random}: the reason depends on other columns you have (older people skip
    a question more often). Leaving rows out can bias the result; taking that other column into
    account helps.
  \item \textbf{Missing not at random}: the reason depends on the missing value itself (very ill
    patients drop out; concentrations below the detection limit). No simple fix removes the bias;
    say so in your report.
\end{itemize}

\textbf{Removing} rows with empty cells (\code{table.remove\_rows\_with\_empty\_cells()}) is
simple and honest when few rows are affected. \textbf{Filling} them in is called
\textbf{imputation}: \code{table.fill\_empty\_cells("marks", with\_the="average")} (also
\code{"median"}, \code{"value above"}, \code{"value below"}, or \code{with\_value=0}). Filling with
the average keeps the average the same, but makes the values look \emph{less spread out} than they
really are, so standard deviations, standard errors and p-values become too small. Filling with the
value above or below suits a series in time order. Whatever you choose, count the empty cells first
(\code{table.count\_empty\_cells()}) and say in your methods how you handled them.

\ideachapter{Which test when?}{ideas:which-test}

First draw the data: a chart shows outliers, skew and odd values that a single p-value hides. Then
answer four questions:

\begin{enumerate}
  \item \textbf{What is the question?} A difference between groups, a change within the same
    rows, a link between two columns, or the shape of a relation?
  \item \textbf{What kind of data?} Numbers (marks, heights, concentrations) or categories (pass or
    fail, species, site)?
  \item \textbf{How many groups, and are they the same rows?} Two groups or more; independent
    groups (different children) or paired measurements (the same children twice)?
  \item \textbf{Are the groups small and the values far from bell-shaped?} Then use the rank-based
    version: \code{based\_on="ranks"}. simplydata checks each group for a bell shape
    (Shapiro--Wilk test) and adds a note when it matters.
\end{enumerate}

\begin{center}\footnotesize
\renewcommand{\arraystretch}{1.25}
\begin{tabularx}{\textwidth}{@{}>{\raggedright\arraybackslash}p{0.25\textwidth}>{\raggedright\arraybackslash}p{0.27\textwidth}>{\raggedright\arraybackslash}X>{\raggedright\arraybackslash}X@{}}
\toprule
\textbf{Question} & \textbf{simplydata word} & \textbf{Test (the usual one)} & \textbf{With \code{based\_on="ranks"}}\\
\midrule
Is one average different from a number you choose? & \code{compare\_average\_to("marks", 60)} & one-sample $t$-test & Wilcoxon signed-rank test\\
Do \textbf{two} groups differ in their average? & \code{compare\_averages("marks", between="group")} & Welch's $t$-test & Mann--Whitney $U$ test\\
Do \textbf{three or more} groups differ? & \code{compare\_averages("marks", between="class")} & Welch's one-way ANOVA & Kruskal--Wallis test\\
\emph{Which} groups differ? & \code{compare\_each\_pair("marks", between="class")} & pairwise Welch $t$-tests, Holm's correction & pairwise Mann--Whitney tests, Holm's correction\\
Did the values change, in the \textbf{same} rows measured twice? & \code{compare\_before\_after("before", "after")} & paired $t$-test & Wilcoxon signed-rank test\\
Do \textbf{counts} of categories differ between groups? & \code{compare\_counts("result", between="class")} & chi-square test of independence & --- (for a $2\times2$ table with small counts: \code{test="fisher"}, Fisher's exact test)\\
Do two columns of numbers go up and down together? & \code{correlation\_between("hours", "marks")} & Pearson's $r$ & \code{method="spearman"}: Spearman's $\rho$\\
How much does $y$ change with $x$? & \code{fit\_straight\_line(x="hours", y="marks")} & linear regression (least squares) & ---\\
Does the data follow a formula (growth, decay, saturation)? & \code{fit\_formula("y = a * exp(b * x)")} & nonlinear least squares & ---\\
\bottomrule
\end{tabularx}
\end{center}

\begin{goodtoknow}
\begin{itemize}
  \item \textbf{Welch, not Student.} For averages simplydata uses Welch's versions, which do not
    assume that the groups have the same spread. They are R's defaults and lose almost nothing when
    the spreads are equal. An expert can still choose: \code{test="student"} or \code{test="anova"}.
  \item \textbf{Paired or not?} If each row of one measurement belongs to a row of the other (the same
    child, the same site), the test must be paired: \code{compare\_before\_after}. Treating paired
    data as independent groups throws away information.
  \item \textbf{Counts, not percentages.} \code{compare\_counts} needs the rows themselves (one row per
    observation), not percentages worked out by hand.
  \item \textbf{One test per question.} Decide the question and the test before looking at the
    results; trying several tests and keeping the best p-value makes false alarms much more
    likely.
  \item Each test gives back a one-row table with every number in it:
    \code{result = table.compare\_averages(...)}, then \code{result.save("test.csv")}.
\end{itemize}
\end{goodtoknow}

\unsection{Writing it up}{ideas:report}

A good results sentence lets a reader check your work. Give:
\begin{enumerate}
  \item the summary for each group (average and SD, or median and IQR) and the number of values $n$;
  \item the effect size, ideally with its 95\,\% confidence interval;
  \item the name of the test and the exact p-value;
  \item in figure captions, what the bars or bands show (SD, SE or 95\,\% CI) and $n$.
\end{enumerate}

For example: ``Class~A scored on average 13.2 marks more than class~C (Welch's $t$-test with Holm's
correction, $p = 0.023$; $n = 18$ per class).'' simplydata's printed answers name each test, and
\code{table.steps()} lists every change you made to the table, in order, for your methods.
).
% Plain words first, then the formal name, so a learner can follow a paper's Methods section.

\unchapter{Describing numbers}{ideas:describe}

Before any test or model, a few numbers describe a column well: where its middle is, and how
spread out it is. simplydata's summary words say each one in a sentence.

\unsection{Average or median?}{ideas:average-median}

The \textbf{average} (formally the \textbf{arithmetic mean}, written $\bar{x}$) adds all the values
and divides by how many there are:
\[
  \bar{x} = \frac{x_1 + x_2 + \dots + x_n}{n}.
\]
The \textbf{median} is the middle value when the values are put in order (with an even number of
values, the average of the two in the middle). Half the values are below it and half above.

Both describe the ``middle'', but they react differently to values far from the rest
(\textbf{outliers}). Five children earn pocket money of 2, 3, 3, 4 and 50 a week: the average is
12.4, though four of the five get 4 or less; the median is 3. One extreme value pulls the average
towards it but hardly moves the median. So:

\begin{itemize}
  \item Use the \textbf{average} when the values are roughly symmetric (a bell shape), and when you
    want to add things up (total = average $\times$ count).
  \item Use the \textbf{median} when the values are lopsided (\textbf{skewed}), such as incomes,
    house prices, waiting times or bacterial counts, or when there are outliers you cannot
    remove.
  \item When the two are far apart, say so: it tells the reader the values are skewed.
\end{itemize}

In simplydata: \code{table.average\_of("marks")}, \code{table.median\_of("marks")},
\code{nums.average()}, \code{nums.median()}, and on charts \code{show="average"} or
\code{show="median"}. Lesson~\ref{lesson:5} draws both.

\unsection{Spread: range, percentiles and standard deviation}{ideas:spread}

Two classes can have the same average mark, one with every child near 65 and the other with
marks from 20 to 100. The \textbf{spread} tells them apart.

\begin{itemize}
  \item The \textbf{range} is the biggest value minus the smallest. It is easy to understand but
    depends on just two values, often the strangest ones.
  \item A \textbf{percentile} is the value that a given share of the values is below:
    \code{nums.percentile(90)} is the value that 90\,\% are below. The median is the 50th
    percentile. The middle half of the values lies between the 25th and 75th percentiles; that
    distance is the \textbf{interquartile range (IQR)}, and it is the box in a box plot
    (\code{chart.boxes}).
  \item The \textbf{standard deviation} (SD, written $s$) is about how far a typical value is from
    the average:
    \[
      s = \sqrt{\frac{\sum_{i=1}^{n} (x_i - \bar{x})^2}{n - 1}}.
    \]
    Its square, $s^2$, is the \textbf{variance}. Dividing by $n-1$ instead of $n$ (Bessel's
    correction) makes the sample variance an unbiased estimate of the variance of everything the
    sample came from; simplydata always divides by $n-1$, as pandas and R do (NumPy divides by $n$
    unless told otherwise).
\end{itemize}

The SD goes with the average and the IQR with the median. For all the marks of the practice table:
``average 64.9 (SD 15.0)'' or ``median 68.5 (IQR 54.3 to 74.8)''.

\unsection{The shape of the values: distributions}{ideas:distribution}

The way values are spread out is their \textbf{distribution}; a histogram
(\code{chart.histogram}) or a density curve (\code{chart.density}) shows it. Many measurements,
such as heights or measuring errors, follow a bell shape called the \textbf{normal} (or
Gaussian) \textbf{distribution}. For a normal distribution about 68\,\% of the values lie within
one SD of the average and about 95\,\% within two SD.

Not everything is bell-shaped. Counts and concentrations are often skewed to the right (many
small values, a few large ones); a log scale (\code{chart.log\_scale\_y()}) often makes them look
more symmetric. Look at the distribution before choosing what to report and which test to use.

\ideachapter{How sure are we?}{ideas:uncertainty}

We almost never measure everything. We measure a \textbf{sample} (18 children, 40 water samples)
and want to learn about the \textbf{population} it came from (all children of that age, the whole
fjord). A number worked out from a sample, such as its average, is an \textbf{estimate}, and
another sample would give a slightly different one. These ideas say how much.

\unsection{Standard deviation or standard error?}{ideas:sd-se}

They sound alike and are often mixed up, but they answer different questions.

\begin{center}\small
\begin{tabularx}{\textwidth}{@{}>{\raggedright\arraybackslash}p{0.19\textwidth}>{\raggedright\arraybackslash}X>{\raggedright\arraybackslash}X@{}}
\toprule
 & \textbf{Standard deviation (SD)} & \textbf{Standard error of the mean (SE, SEM)}\\
\midrule
Describes & the \emph{values}: how spread out they are & the \emph{average}: how precisely the sample pins it down\\
Formula & $s$ (above) & $\mathrm{SE} = s / \sqrt{n}$\\
More data & stays about the same & gets smaller (4 times the data halves it)\\
Use it to & show the spread of individuals & compare averages, build confidence intervals\\
simplydata & \code{error="standard deviation"} & \code{error="standard error"}\\
\bottomrule
\end{tabularx}
\end{center}

Error bars of the SE are always shorter than those of the SD, and can look reassuringly small. That
is why a figure caption must say which one was drawn; simplydata prints a line saying it, ready
for your caption.

\unsection{Confidence intervals}{ideas:ci}

A \textbf{95\,\% confidence interval} (CI) is a range of plausible values for the true number in
the population, worked out from the sample. For an average it is
\[
  \bar{x} \pm t_{0.975,\,n-1} \times \mathrm{SE},
\]
where $t_{0.975,\,n-1}$ is a number from the $t$ distribution: about 2 for most samples, a bit
more for small ones (2.26 for $n = 10$, 2.09 for $n = 20$).

What ``95\,\%'' means, exactly: if you repeated the whole study many times and built an interval
each time, about 95 out of 100 of those intervals would contain the true value. It does
\emph{not} mean that the true value moves around, nor that 95\,\% of the individual values lie
inside (that would be about $\bar{x} \pm 2\,\mathrm{SD}$ for a bell shape).

\begin{itemize}
  \item A \textbf{narrow} interval means the data pin the value down well; a \textbf{wide} one
    means you need more data.
  \item An interval says more than a p-value: it gives the size of the effect \emph{and} how sure
    you are of it. \code{table.fit\_straight\_line} prints the slope with its 95\,\% CI.
  \item On charts: \code{error="confidence interval"} on bars and lines, and
    \code{fit\_line(confidence\_band=True)} for a fitted line.
\end{itemize}

\unsection{Why averages behave well: the central limit theorem}{ideas:clt}

Even when the values themselves are skewed, the \emph{average} of many of them is close to
bell-shaped; this is the \textbf{central limit theorem}. It is why tests on averages (t-tests,
ANOVA) usually work well with 30 or more values per group, even for values that are not
bell-shaped. With small groups of skewed values, rank-based tests are safer
(\code{based\_on="ranks"}).

\ideachapter{Testing: real difference, or chance?}{ideas:testing}

\unsection{The null hypothesis and the p-value}{ideas:p-value}

Two groups almost never have exactly the same average, even when nothing real separates them:
samples differ by chance. A statistical test asks whether a difference is bigger than chance alone
would often produce. It works like a court case, where the accused is innocent until proven
guilty:

\begin{enumerate}
  \item Assume there is \textbf{no real difference}. This assumption is the \textbf{null
    hypothesis} ($H_0$). The \textbf{alternative hypothesis} ($H_1$) is that there is one.
  \item Work out how surprising your data would be if $H_0$ were true. The \textbf{p-value} is the
    probability of a difference at least as big as the one you saw, if there were really no
    difference.
  \item A small p-value means the data are hard to explain by chance alone. By custom, below
    $0.05$ is called \textbf{statistically significant}; that threshold is the
    \textbf{significance level} $\alpha$.
\end{enumerate}

simplydata says the p-value in words: below 0.001 ``very unlikely to be chance'', below 0.05
``unlikely to be chance'', otherwise ``could easily be chance''.

\begin{goodtoknow}
What a p-value is \textbf{not}:
\begin{itemize}
  \item It is \textbf{not} the probability that $H_0$ is true, nor the probability that the result
    is a fluke.
  \item A small p-value does \textbf{not} say the difference is big or important. With thousands
    of values, tiny, meaningless differences become ``significant''. The \emph{effect size} says
    how big it is.
  \item A large p-value does \textbf{not} show that there is no difference. It only says these
    data cannot tell. ``Absence of evidence is not evidence of absence.''
  \item 0.049 and 0.051 are practically the same evidence. Report the exact p-value, not just
    ``$p < 0.05$''.
\end{itemize}
\end{goodtoknow}

\unsection{Two kinds of mistake, and power}{ideas:errors}

\begin{itemize}
  \item A \textbf{false alarm} (\textbf{type I error}, a false positive): the test says there is a
    difference when there is none. With $\alpha = 0.05$ this happens in about 1 test in 20 when
    $H_0$ is true.
  \item A \textbf{missed difference} (\textbf{type II error}, a false negative): there is a real
    difference, but the test does not find it. Small samples miss differences often.
\end{itemize}

The \textbf{power} of a study is the chance of finding a difference of a given size if it is
really there. More data, less noisy measurements and bigger true differences all raise power. A
study with low power can only find large effects; plan the sample size before collecting data.

\unsection{Many tests at once}{ideas:multiple}

Every test at $\alpha = 0.05$ carries a 5\,\% risk of a false alarm, and the risks add up: with 20
tests on pure noise, you expect about one ``significant'' result. Correcting for
\textbf{multiple comparisons} keeps the chance of \emph{any} false alarm at 5\,\%.
\code{table.compare\_each\_pair} uses \textbf{Holm's method}, which is never less powerful than the
simpler Bonferroni correction (multiply each p-value by the number of tests). Decide which
comparisons you will make before you look at the data.

\unsection{Effect sizes: how big is the difference?}{ideas:effect-size}

The p-value answers ``could it be chance?''. The \textbf{effect size} answers ``how big is it?'',
which is usually what matters. Report both, ideally with a confidence interval for the effect.

\begin{center}\small
\begin{tabularx}{\textwidth}{@{}>{\raggedright\arraybackslash}p{0.26\textwidth}>{\raggedright\arraybackslash}X>{\raggedright\arraybackslash}p{0.25\textwidth}@{}}
\toprule
\textbf{Effect size} & \textbf{What it says} & \textbf{Given by}\\
\midrule
Difference in averages & in the units of the data (``7.6 marks higher''): often the clearest & \code{compare\_averages}, \code{compare\_each\_pair}\\
Cohen's $d$, Hedges' $g$ & the difference in standard deviations; $g$ corrects $d$ for small samples. Rough guide: 0.2 small, 0.5 medium, 0.8 large & \code{compare\_averages} (2 groups), \code{compare\_average\_to}, \code{compare\_before\_after} ($d_z$)\\
Omega-squared $\omega^2$ & the share of all the differences that the groups explain (0 to 1) & \code{compare\_averages} (3+ groups)\\
Cramér's $V$ & how strongly two categories are linked, 0 (none) to 1 & \code{compare\_counts}\\
Correlation $r$ & how closely two columns go up and down together, $-1$ to $1$ & \code{correlation\_between}\\
$R^2$ & the share of the differences in $y$ that a line or curve explains & \code{fit\_straight\_line}, \code{fit\_formula}\\
\bottomrule
\end{tabularx}
\end{center}

The ``small, medium, large'' labels (from Jacob Cohen) are only a last resort: a ``small''
difference in survival can be important, and a ``large'' one in a lab assay may not be.

\ideachapter{Links between columns}{ideas:links}

\unsection{Correlation}{ideas:correlation}

A \textbf{correlation} measures how closely two columns of numbers go up and down together, from
$-1$ (one always goes down when the other goes up) through $0$ (no straight-line link) to $1$.

\begin{itemize}
  \item \textbf{Pearson's $r$} (the usual one) measures a \emph{straight-line} link. It is
    sensitive to outliers, and can be near 0 for a strong curved link (a U shape).
  \item \textbf{Spearman's $\rho$} (rho) uses ranks: it measures any link that keeps going one way
    (always up, or always down), and is robust to outliers.
    \code{table.correlation\_between("hours", "marks", method="spearman")}.
  \item $r^2$ is the share of the differences in one column that a straight line through the other
    explains. $r = 0.5$ explains only a quarter.
\end{itemize}

Always draw the dots (\code{chart.dots}) before trusting a correlation: very different patterns
can have the same $r$ (Anscombe's quartet is the famous example).

\unsection{Correlation is not cause}{ideas:cause}

Ice cream sales and drownings go up together, but ice cream does not cause drowning: hot weather
drives both. When two things go together there are several possible reasons:

\begin{enumerate}
  \item A causes B.
  \item B causes A (\textbf{reverse causation}): do people exercise because they are healthy, or
    are they healthy because they exercise?
  \item A third thing causes both (a \textbf{confounder}): hot weather, age, wealth, or the size of
    a hospital.
  \item Chance: test enough pairs of columns and some will line up by luck.
\end{enumerate}

Only an \textbf{experiment} where you change A yourself, ideally by random assignment (a
randomised controlled trial), shows cause directly. With observational data, say ``is associated
with'', not ``causes''. simplydata's link words remind you: ``A link does not show that one thing
causes the other.''

\unsection{Fitting a line or a formula}{ideas:regression}

\textbf{Linear regression} finds the straight line $y = a + b\,x$ that goes as close as possible to
the dots, by making the sum of the squared vertical distances from the dots to the line as small as
possible (\textbf{least squares}). The \textbf{slope} $b$ says how much $y$ changes for each 1 more
of $x$; its p-value tests whether the slope could be 0, and its confidence interval gives a range of
plausible slopes. $R^2$ says how much of the differences in $y$ the line explains.
\code{table.fit\_straight\_line(x="hours", y="marks")} prints all of these;
\code{chart.fit\_line()} draws the line.

When the relation is curved, such as growth that multiplies by the same share each year, write the
formula yourself: \code{table.fit\_formula("population = a * exp(b * (year - 2000))")} finds $a$ and
$b$ by \textbf{nonlinear least squares}. A fitted line or curve describes the data you have;
predicting far outside their range (\textbf{extrapolation}) is risky.

\ideachapter{The shape of a table}{ideas:tables}

\unsection{Tidy data: long and wide tables}{ideas:tidy}

The same marks can be stored in two shapes:

\begin{center}\small
\begin{tabular}{@{}lll@{\hspace{3em}}llll@{}}
\multicolumn{3}{@{}l}{\textbf{Long (tidy)}} & \multicolumn{4}{l@{}}{\textbf{Wide}}\\
\toprule
name & subject & marks & name & Maths & Science & English\\
\midrule
Asha & Maths & 61 & Asha & 61 & 66 & 77\\
Asha & Science & 66 & Ben & 73 & 98 & 71\\
Asha & English & 77 & & & & \\
Ben & Maths & 73 & & & & \\
\multicolumn{3}{@{}l}{\dots} & & & & \\
\bottomrule
\end{tabular}
\end{center}

A table is \textbf{tidy} (Hadley Wickham's term) when each \textbf{variable} is a column, each
\textbf{observation} is a row, and each value has its own cell. The long table is tidy: it has one
column for what was measured (\code{marks}) and columns that say which child and which subject.

\begin{itemize}
  \item \textbf{Long tables} suit charts, groups and tests: \code{chart.bars(x="subject", y="marks")}
    and \code{compare\_averages("marks", between="subject")} need the subject as a column.
  \item \textbf{Wide tables} suit sums across columns (a total per child) and paired tests
    (\code{compare\_before\_after("Maths", "Science")}), and are how people often type data into
    a spreadsheet.
  \item \code{table.turn\_columns\_into\_rows} makes a wide table long (in pandas: \code{melt}; in R:
    \code{pivot\_longer}); \code{table.turn\_rows\_into\_columns} makes a long table wide (\code{pivot}).
\end{itemize}

\unsection{Missing values and imputation}{ideas:missing}

Real data has empty cells: a sample was lost, a question skipped. What to do depends on
\emph{why} they are empty:

\begin{itemize}
  \item \textbf{Missing completely at random}: the reason has nothing to do with the data (a tube
    dropped). Leaving those rows out loses some power but does not bias the result.
  \item \textbf{Missing at random}: the reason depends on other columns you have (older people skip
    a question more often). Leaving rows out can bias the result; taking that other column into
    account helps.
  \item \textbf{Missing not at random}: the reason depends on the missing value itself (very ill
    patients drop out; concentrations below the detection limit). No simple fix removes the bias;
    say so in your report.
\end{itemize}

\textbf{Removing} rows with empty cells (\code{table.remove\_rows\_with\_empty\_cells()}) is
simple and honest when few rows are affected. \textbf{Filling} them in is called
\textbf{imputation}: \code{table.fill\_empty\_cells("marks", with\_the="average")} (also
\code{"median"}, \code{"value above"}, \code{"value below"}, or \code{with\_value=0}). Filling with
the average keeps the average the same, but makes the values look \emph{less spread out} than they
really are, so standard deviations, standard errors and p-values become too small. Filling with the
value above or below suits a series in time order. Whatever you choose, count the empty cells first
(\code{table.count\_empty\_cells()}) and say in your methods how you handled them.

\ideachapter{Which test when?}{ideas:which-test}

First draw the data: a chart shows outliers, skew and odd values that a single p-value hides. Then
answer four questions:

\begin{enumerate}
  \item \textbf{What is the question?} A difference between groups, a change within the same
    rows, a link between two columns, or the shape of a relation?
  \item \textbf{What kind of data?} Numbers (marks, heights, concentrations) or categories (pass or
    fail, species, site)?
  \item \textbf{How many groups, and are they the same rows?} Two groups or more; independent
    groups (different children) or paired measurements (the same children twice)?
  \item \textbf{Are the groups small and the values far from bell-shaped?} Then use the rank-based
    version: \code{based\_on="ranks"}. simplydata checks each group for a bell shape
    (Shapiro--Wilk test) and adds a note when it matters.
\end{enumerate}

\begin{center}\footnotesize
\renewcommand{\arraystretch}{1.25}
\begin{tabularx}{\textwidth}{@{}>{\raggedright\arraybackslash}p{0.25\textwidth}>{\raggedright\arraybackslash}p{0.27\textwidth}>{\raggedright\arraybackslash}X>{\raggedright\arraybackslash}X@{}}
\toprule
\textbf{Question} & \textbf{simplydata word} & \textbf{Test (the usual one)} & \textbf{With \code{based\_on="ranks"}}\\
\midrule
Is one average different from a number you choose? & \code{compare\_average\_to("marks", 60)} & one-sample $t$-test & Wilcoxon signed-rank test\\
Do \textbf{two} groups differ in their average? & \code{compare\_averages("marks", between="group")} & Welch's $t$-test & Mann--Whitney $U$ test\\
Do \textbf{three or more} groups differ? & \code{compare\_averages("marks", between="class")} & Welch's one-way ANOVA & Kruskal--Wallis test\\
\emph{Which} groups differ? & \code{compare\_each\_pair("marks", between="class")} & pairwise Welch $t$-tests, Holm's correction & pairwise Mann--Whitney tests, Holm's correction\\
Did the values change, in the \textbf{same} rows measured twice? & \code{compare\_before\_after("before", "after")} & paired $t$-test & Wilcoxon signed-rank test\\
Do \textbf{counts} of categories differ between groups? & \code{compare\_counts("result", between="class")} & chi-square test of independence & --- (for a $2\times2$ table with small counts: \code{test="fisher"}, Fisher's exact test)\\
Do two columns of numbers go up and down together? & \code{correlation\_between("hours", "marks")} & Pearson's $r$ & \code{method="spearman"}: Spearman's $\rho$\\
How much does $y$ change with $x$? & \code{fit\_straight\_line(x="hours", y="marks")} & linear regression (least squares) & ---\\
Does the data follow a formula (growth, decay, saturation)? & \code{fit\_formula("y = a * exp(b * x)")} & nonlinear least squares & ---\\
\bottomrule
\end{tabularx}
\end{center}

\begin{goodtoknow}
\begin{itemize}
  \item \textbf{Welch, not Student.} For averages simplydata uses Welch's versions, which do not
    assume that the groups have the same spread. They are R's defaults and lose almost nothing when
    the spreads are equal. An expert can still choose: \code{test="student"} or \code{test="anova"}.
  \item \textbf{Paired or not?} If each row of one measurement belongs to a row of the other (the same
    child, the same site), the test must be paired: \code{compare\_before\_after}. Treating paired
    data as independent groups throws away information.
  \item \textbf{Counts, not percentages.} \code{compare\_counts} needs the rows themselves (one row per
    observation), not percentages worked out by hand.
  \item \textbf{One test per question.} Decide the question and the test before looking at the
    results; trying several tests and keeping the best p-value makes false alarms much more
    likely.
  \item Each test gives back a one-row table with every number in it:
    \code{result = table.compare\_averages(...)}, then \code{result.save("test.csv")}.
\end{itemize}
\end{goodtoknow}

\unsection{Writing it up}{ideas:report}

A good results sentence lets a reader check your work. Give:
\begin{enumerate}
  \item the summary for each group (average and SD, or median and IQR) and the number of values $n$;
  \item the effect size, ideally with its 95\,\% confidence interval;
  \item the name of the test and the exact p-value;
  \item in figure captions, what the bars or bands show (SD, SE or 95\,\% CI) and $n$.
\end{enumerate}

For example: ``Class~A scored on average 13.2 marks more than class~C (Welch's $t$-test with Holm's
correction, $p = 0.023$; $n = 18$ per class).'' simplydata's printed answers name each test, and
\code{table.steps()} lists every change you made to the table, in order, for your methods.
).
% Plain words first, then the formal name, so a learner can follow a paper's Methods section.

\unchapter{Describing numbers}{ideas:describe}

Before any test or model, a few numbers describe a column well: where its middle is, and how
spread out it is. simplydata's summary words say each one in a sentence.

\unsection{Average or median?}{ideas:average-median}

The \textbf{average} (formally the \textbf{arithmetic mean}, written $\bar{x}$) adds all the values
and divides by how many there are:
\[
  \bar{x} = \frac{x_1 + x_2 + \dots + x_n}{n}.
\]
The \textbf{median} is the middle value when the values are put in order (with an even number of
values, the average of the two in the middle). Half the values are below it and half above.

Both describe the ``middle'', but they react differently to values far from the rest
(\textbf{outliers}). Five children earn pocket money of 2, 3, 3, 4 and 50 a week: the average is
12.4, though four of the five get 4 or less; the median is 3. One extreme value pulls the average
towards it but hardly moves the median. So:

\begin{itemize}
  \item Use the \textbf{average} when the values are roughly symmetric (a bell shape), and when you
    want to add things up (total = average $\times$ count).
  \item Use the \textbf{median} when the values are lopsided (\textbf{skewed}), such as incomes,
    house prices, waiting times or bacterial counts, or when there are outliers you cannot
    remove.
  \item When the two are far apart, say so: it tells the reader the values are skewed.
\end{itemize}

In simplydata: \code{table.average\_of("marks")}, \code{table.median\_of("marks")},
\code{nums.average()}, \code{nums.median()}, and on charts \code{show="average"} or
\code{show="median"}. Lesson~\ref{lesson:5} draws both.

\unsection{Spread: range, percentiles and standard deviation}{ideas:spread}

Two classes can have the same average mark, one with every child near 65 and the other with
marks from 20 to 100. The \textbf{spread} tells them apart.

\begin{itemize}
  \item The \textbf{range} is the biggest value minus the smallest. It is easy to understand but
    depends on just two values, often the strangest ones.
  \item A \textbf{percentile} is the value that a given share of the values is below:
    \code{nums.percentile(90)} is the value that 90\,\% are below. The median is the 50th
    percentile. The middle half of the values lies between the 25th and 75th percentiles; that
    distance is the \textbf{interquartile range (IQR)}, and it is the box in a box plot
    (\code{chart.boxes}).
  \item The \textbf{standard deviation} (SD, written $s$) is about how far a typical value is from
    the average:
    \[
      s = \sqrt{\frac{\sum_{i=1}^{n} (x_i - \bar{x})^2}{n - 1}}.
    \]
    Its square, $s^2$, is the \textbf{variance}. Dividing by $n-1$ instead of $n$ (Bessel's
    correction) makes the sample variance an unbiased estimate of the variance of everything the
    sample came from; simplydata always divides by $n-1$, as pandas and R do (NumPy divides by $n$
    unless told otherwise).
\end{itemize}

The SD goes with the average and the IQR with the median. For all the marks of the practice table:
``average 64.9 (SD 15.0)'' or ``median 68.5 (IQR 54.3 to 74.8)''.

\unsection{The shape of the values: distributions}{ideas:distribution}

The way values are spread out is their \textbf{distribution}; a histogram
(\code{chart.histogram}) or a density curve (\code{chart.density}) shows it. Many measurements,
such as heights or measuring errors, follow a bell shape called the \textbf{normal} (or
Gaussian) \textbf{distribution}. For a normal distribution about 68\,\% of the values lie within
one SD of the average and about 95\,\% within two SD.

Not everything is bell-shaped. Counts and concentrations are often skewed to the right (many
small values, a few large ones); a log scale (\code{chart.log\_scale\_y()}) often makes them look
more symmetric. Look at the distribution before choosing what to report and which test to use.

\ideachapter{How sure are we?}{ideas:uncertainty}

We almost never measure everything. We measure a \textbf{sample} (18 children, 40 water samples)
and want to learn about the \textbf{population} it came from (all children of that age, the whole
fjord). A number worked out from a sample, such as its average, is an \textbf{estimate}, and
another sample would give a slightly different one. These ideas say how much.

\unsection{Standard deviation or standard error?}{ideas:sd-se}

They sound alike and are often mixed up, but they answer different questions.

\begin{center}\small
\begin{tabularx}{\textwidth}{@{}>{\raggedright\arraybackslash}p{0.19\textwidth}>{\raggedright\arraybackslash}X>{\raggedright\arraybackslash}X@{}}
\toprule
 & \textbf{Standard deviation (SD)} & \textbf{Standard error of the mean (SE, SEM)}\\
\midrule
Describes & the \emph{values}: how spread out they are & the \emph{average}: how precisely the sample pins it down\\
Formula & $s$ (above) & $\mathrm{SE} = s / \sqrt{n}$\\
More data & stays about the same & gets smaller (4 times the data halves it)\\
Use it to & show the spread of individuals & compare averages, build confidence intervals\\
simplydata & \code{error="standard deviation"} & \code{error="standard error"}\\
\bottomrule
\end{tabularx}
\end{center}

Error bars of the SE are always shorter than those of the SD, and can look reassuringly small. That
is why a figure caption must say which one was drawn; simplydata prints a line saying it, ready
for your caption.

\unsection{Confidence intervals}{ideas:ci}

A \textbf{95\,\% confidence interval} (CI) is a range of plausible values for the true number in
the population, worked out from the sample. For an average it is
\[
  \bar{x} \pm t_{0.975,\,n-1} \times \mathrm{SE},
\]
where $t_{0.975,\,n-1}$ is a number from the $t$ distribution: about 2 for most samples, a bit
more for small ones (2.26 for $n = 10$, 2.09 for $n = 20$).

What ``95\,\%'' means, exactly: if you repeated the whole study many times and built an interval
each time, about 95 out of 100 of those intervals would contain the true value. It does
\emph{not} mean that the true value moves around, nor that 95\,\% of the individual values lie
inside (that would be about $\bar{x} \pm 2\,\mathrm{SD}$ for a bell shape).

\begin{itemize}
  \item A \textbf{narrow} interval means the data pin the value down well; a \textbf{wide} one
    means you need more data.
  \item An interval says more than a p-value: it gives the size of the effect \emph{and} how sure
    you are of it. \code{table.fit\_straight\_line} prints the slope with its 95\,\% CI.
  \item On charts: \code{error="confidence interval"} on bars and lines, and
    \code{fit\_line(confidence\_band=True)} for a fitted line.
\end{itemize}

\unsection{Why averages behave well: the central limit theorem}{ideas:clt}

Even when the values themselves are skewed, the \emph{average} of many of them is close to
bell-shaped; this is the \textbf{central limit theorem}. It is why tests on averages (t-tests,
ANOVA) usually work well with 30 or more values per group, even for values that are not
bell-shaped. With small groups of skewed values, rank-based tests are safer
(\code{based\_on="ranks"}).

\ideachapter{Testing: real difference, or chance?}{ideas:testing}

\unsection{The null hypothesis and the p-value}{ideas:p-value}

Two groups almost never have exactly the same average, even when nothing real separates them:
samples differ by chance. A statistical test asks whether a difference is bigger than chance alone
would often produce. It works like a court case, where the accused is innocent until proven
guilty:

\begin{enumerate}
  \item Assume there is \textbf{no real difference}. This assumption is the \textbf{null
    hypothesis} ($H_0$). The \textbf{alternative hypothesis} ($H_1$) is that there is one.
  \item Work out how surprising your data would be if $H_0$ were true. The \textbf{p-value} is the
    probability of a difference at least as big as the one you saw, if there were really no
    difference.
  \item A small p-value means the data are hard to explain by chance alone. By custom, below
    $0.05$ is called \textbf{statistically significant}; that threshold is the
    \textbf{significance level} $\alpha$.
\end{enumerate}

simplydata says the p-value in words: below 0.001 ``very unlikely to be chance'', below 0.05
``unlikely to be chance'', otherwise ``could easily be chance''.

\begin{goodtoknow}
What a p-value is \textbf{not}:
\begin{itemize}
  \item It is \textbf{not} the probability that $H_0$ is true, nor the probability that the result
    is a fluke.
  \item A small p-value does \textbf{not} say the difference is big or important. With thousands
    of values, tiny, meaningless differences become ``significant''. The \emph{effect size} says
    how big it is.
  \item A large p-value does \textbf{not} show that there is no difference. It only says these
    data cannot tell. ``Absence of evidence is not evidence of absence.''
  \item 0.049 and 0.051 are practically the same evidence. Report the exact p-value, not just
    ``$p < 0.05$''.
\end{itemize}
\end{goodtoknow}

\unsection{Two kinds of mistake, and power}{ideas:errors}

\begin{itemize}
  \item A \textbf{false alarm} (\textbf{type I error}, a false positive): the test says there is a
    difference when there is none. With $\alpha = 0.05$ this happens in about 1 test in 20 when
    $H_0$ is true.
  \item A \textbf{missed difference} (\textbf{type II error}, a false negative): there is a real
    difference, but the test does not find it. Small samples miss differences often.
\end{itemize}

The \textbf{power} of a study is the chance of finding a difference of a given size if it is
really there. More data, less noisy measurements and bigger true differences all raise power. A
study with low power can only find large effects; plan the sample size before collecting data.

\unsection{Many tests at once}{ideas:multiple}

Every test at $\alpha = 0.05$ carries a 5\,\% risk of a false alarm, and the risks add up: with 20
tests on pure noise, you expect about one ``significant'' result. Correcting for
\textbf{multiple comparisons} keeps the chance of \emph{any} false alarm at 5\,\%.
\code{table.compare\_each\_pair} uses \textbf{Holm's method}, which is never less powerful than the
simpler Bonferroni correction (multiply each p-value by the number of tests). Decide which
comparisons you will make before you look at the data.

\unsection{Effect sizes: how big is the difference?}{ideas:effect-size}

The p-value answers ``could it be chance?''. The \textbf{effect size} answers ``how big is it?'',
which is usually what matters. Report both, ideally with a confidence interval for the effect.

\begin{center}\small
\begin{tabularx}{\textwidth}{@{}>{\raggedright\arraybackslash}p{0.26\textwidth}>{\raggedright\arraybackslash}X>{\raggedright\arraybackslash}p{0.25\textwidth}@{}}
\toprule
\textbf{Effect size} & \textbf{What it says} & \textbf{Given by}\\
\midrule
Difference in averages & in the units of the data (``7.6 marks higher''): often the clearest & \code{compare\_averages}, \code{compare\_each\_pair}\\
Cohen's $d$, Hedges' $g$ & the difference in standard deviations; $g$ corrects $d$ for small samples. Rough guide: 0.2 small, 0.5 medium, 0.8 large & \code{compare\_averages} (2 groups), \code{compare\_average\_to}, \code{compare\_before\_after} ($d_z$)\\
Omega-squared $\omega^2$ & the share of all the differences that the groups explain (0 to 1) & \code{compare\_averages} (3+ groups)\\
Cramér's $V$ & how strongly two categories are linked, 0 (none) to 1 & \code{compare\_counts}\\
Correlation $r$ & how closely two columns go up and down together, $-1$ to $1$ & \code{correlation\_between}\\
$R^2$ & the share of the differences in $y$ that a line or curve explains & \code{fit\_straight\_line}, \code{fit\_formula}\\
\bottomrule
\end{tabularx}
\end{center}

The ``small, medium, large'' labels (from Jacob Cohen) are only a last resort: a ``small''
difference in survival can be important, and a ``large'' one in a lab assay may not be.

\ideachapter{Links between columns}{ideas:links}

\unsection{Correlation}{ideas:correlation}

A \textbf{correlation} measures how closely two columns of numbers go up and down together, from
$-1$ (one always goes down when the other goes up) through $0$ (no straight-line link) to $1$.

\begin{itemize}
  \item \textbf{Pearson's $r$} (the usual one) measures a \emph{straight-line} link. It is
    sensitive to outliers, and can be near 0 for a strong curved link (a U shape).
  \item \textbf{Spearman's $\rho$} (rho) uses ranks: it measures any link that keeps going one way
    (always up, or always down), and is robust to outliers.
    \code{table.correlation\_between("hours", "marks", method="spearman")}.
  \item $r^2$ is the share of the differences in one column that a straight line through the other
    explains. $r = 0.5$ explains only a quarter.
\end{itemize}

Always draw the dots (\code{chart.dots}) before trusting a correlation: very different patterns
can have the same $r$ (Anscombe's quartet is the famous example).

\unsection{Correlation is not cause}{ideas:cause}

Ice cream sales and drownings go up together, but ice cream does not cause drowning: hot weather
drives both. When two things go together there are several possible reasons:

\begin{enumerate}
  \item A causes B.
  \item B causes A (\textbf{reverse causation}): do people exercise because they are healthy, or
    are they healthy because they exercise?
  \item A third thing causes both (a \textbf{confounder}): hot weather, age, wealth, or the size of
    a hospital.
  \item Chance: test enough pairs of columns and some will line up by luck.
\end{enumerate}

Only an \textbf{experiment} where you change A yourself, ideally by random assignment (a
randomised controlled trial), shows cause directly. With observational data, say ``is associated
with'', not ``causes''. simplydata's link words remind you: ``A link does not show that one thing
causes the other.''

\unsection{Fitting a line or a formula}{ideas:regression}

\textbf{Linear regression} finds the straight line $y = a + b\,x$ that goes as close as possible to
the dots, by making the sum of the squared vertical distances from the dots to the line as small as
possible (\textbf{least squares}). The \textbf{slope} $b$ says how much $y$ changes for each 1 more
of $x$; its p-value tests whether the slope could be 0, and its confidence interval gives a range of
plausible slopes. $R^2$ says how much of the differences in $y$ the line explains.
\code{table.fit\_straight\_line(x="hours", y="marks")} prints all of these;
\code{chart.fit\_line()} draws the line.

When the relation is curved, such as growth that multiplies by the same share each year, write the
formula yourself: \code{table.fit\_formula("population = a * exp(b * (year - 2000))")} finds $a$ and
$b$ by \textbf{nonlinear least squares}. A fitted line or curve describes the data you have;
predicting far outside their range (\textbf{extrapolation}) is risky.

\ideachapter{The shape of a table}{ideas:tables}

\unsection{Tidy data: long and wide tables}{ideas:tidy}

The same marks can be stored in two shapes:

\begin{center}\small
\begin{tabular}{@{}lll@{\hspace{3em}}llll@{}}
\multicolumn{3}{@{}l}{\textbf{Long (tidy)}} & \multicolumn{4}{l@{}}{\textbf{Wide}}\\
\toprule
name & subject & marks & name & Maths & Science & English\\
\midrule
Asha & Maths & 61 & Asha & 61 & 66 & 77\\
Asha & Science & 66 & Ben & 73 & 98 & 71\\
Asha & English & 77 & & & & \\
Ben & Maths & 73 & & & & \\
\multicolumn{3}{@{}l}{\dots} & & & & \\
\bottomrule
\end{tabular}
\end{center}

A table is \textbf{tidy} (Hadley Wickham's term) when each \textbf{variable} is a column, each
\textbf{observation} is a row, and each value has its own cell. The long table is tidy: it has one
column for what was measured (\code{marks}) and columns that say which child and which subject.

\begin{itemize}
  \item \textbf{Long tables} suit charts, groups and tests: \code{chart.bars(x="subject", y="marks")}
    and \code{compare\_averages("marks", between="subject")} need the subject as a column.
  \item \textbf{Wide tables} suit sums across columns (a total per child) and paired tests
    (\code{compare\_before\_after("Maths", "Science")}), and are how people often type data into
    a spreadsheet.
  \item \code{table.turn\_columns\_into\_rows} makes a wide table long (in pandas: \code{melt}; in R:
    \code{pivot\_longer}); \code{table.turn\_rows\_into\_columns} makes a long table wide (\code{pivot}).
\end{itemize}

\unsection{Missing values and imputation}{ideas:missing}

Real data has empty cells: a sample was lost, a question skipped. What to do depends on
\emph{why} they are empty:

\begin{itemize}
  \item \textbf{Missing completely at random}: the reason has nothing to do with the data (a tube
    dropped). Leaving those rows out loses some power but does not bias the result.
  \item \textbf{Missing at random}: the reason depends on other columns you have (older people skip
    a question more often). Leaving rows out can bias the result; taking that other column into
    account helps.
  \item \textbf{Missing not at random}: the reason depends on the missing value itself (very ill
    patients drop out; concentrations below the detection limit). No simple fix removes the bias;
    say so in your report.
\end{itemize}

\textbf{Removing} rows with empty cells (\code{table.remove\_rows\_with\_empty\_cells()}) is
simple and honest when few rows are affected. \textbf{Filling} them in is called
\textbf{imputation}: \code{table.fill\_empty\_cells("marks", with\_the="average")} (also
\code{"median"}, \code{"value above"}, \code{"value below"}, or \code{with\_value=0}). Filling with
the average keeps the average the same, but makes the values look \emph{less spread out} than they
really are, so standard deviations, standard errors and p-values become too small. Filling with the
value above or below suits a series in time order. Whatever you choose, count the empty cells first
(\code{table.count\_empty\_cells()}) and say in your methods how you handled them.

\ideachapter{Which test when?}{ideas:which-test}

First draw the data: a chart shows outliers, skew and odd values that a single p-value hides. Then
answer four questions:

\begin{enumerate}
  \item \textbf{What is the question?} A difference between groups, a change within the same
    rows, a link between two columns, or the shape of a relation?
  \item \textbf{What kind of data?} Numbers (marks, heights, concentrations) or categories (pass or
    fail, species, site)?
  \item \textbf{How many groups, and are they the same rows?} Two groups or more; independent
    groups (different children) or paired measurements (the same children twice)?
  \item \textbf{Are the groups small and the values far from bell-shaped?} Then use the rank-based
    version: \code{based\_on="ranks"}. simplydata checks each group for a bell shape
    (Shapiro--Wilk test) and adds a note when it matters.
\end{enumerate}

\begin{center}\footnotesize
\renewcommand{\arraystretch}{1.25}
\begin{tabularx}{\textwidth}{@{}>{\raggedright\arraybackslash}p{0.25\textwidth}>{\raggedright\arraybackslash}p{0.27\textwidth}>{\raggedright\arraybackslash}X>{\raggedright\arraybackslash}X@{}}
\toprule
\textbf{Question} & \textbf{simplydata word} & \textbf{Test (the usual one)} & \textbf{With \code{based\_on="ranks"}}\\
\midrule
Is one average different from a number you choose? & \code{compare\_average\_to("marks", 60)} & one-sample $t$-test & Wilcoxon signed-rank test\\
Do \textbf{two} groups differ in their average? & \code{compare\_averages("marks", between="group")} & Welch's $t$-test & Mann--Whitney $U$ test\\
Do \textbf{three or more} groups differ? & \code{compare\_averages("marks", between="class")} & Welch's one-way ANOVA & Kruskal--Wallis test\\
\emph{Which} groups differ? & \code{compare\_each\_pair("marks", between="class")} & pairwise Welch $t$-tests, Holm's correction & pairwise Mann--Whitney tests, Holm's correction\\
Did the values change, in the \textbf{same} rows measured twice? & \code{compare\_before\_after("before", "after")} & paired $t$-test & Wilcoxon signed-rank test\\
Do \textbf{counts} of categories differ between groups? & \code{compare\_counts("result", between="class")} & chi-square test of independence & --- (for a $2\times2$ table with small counts: \code{test="fisher"}, Fisher's exact test)\\
Do two columns of numbers go up and down together? & \code{correlation\_between("hours", "marks")} & Pearson's $r$ & \code{method="spearman"}: Spearman's $\rho$\\
How much does $y$ change with $x$? & \code{fit\_straight\_line(x="hours", y="marks")} & linear regression (least squares) & ---\\
Does the data follow a formula (growth, decay, saturation)? & \code{fit\_formula("y = a * exp(b * x)")} & nonlinear least squares & ---\\
\bottomrule
\end{tabularx}
\end{center}

\begin{goodtoknow}
\begin{itemize}
  \item \textbf{Welch, not Student.} For averages simplydata uses Welch's versions, which do not
    assume that the groups have the same spread. They are R's defaults and lose almost nothing when
    the spreads are equal. An expert can still choose: \code{test="student"} or \code{test="anova"}.
  \item \textbf{Paired or not?} If each row of one measurement belongs to a row of the other (the same
    child, the same site), the test must be paired: \code{compare\_before\_after}. Treating paired
    data as independent groups throws away information.
  \item \textbf{Counts, not percentages.} \code{compare\_counts} needs the rows themselves (one row per
    observation), not percentages worked out by hand.
  \item \textbf{One test per question.} Decide the question and the test before looking at the
    results; trying several tests and keeping the best p-value makes false alarms much more
    likely.
  \item Each test gives back a one-row table with every number in it:
    \code{result = table.compare\_averages(...)}, then \code{result.save("test.csv")}.
\end{itemize}
\end{goodtoknow}

\unsection{Writing it up}{ideas:report}

A good results sentence lets a reader check your work. Give:
\begin{enumerate}
  \item the summary for each group (average and SD, or median and IQR) and the number of values $n$;
  \item the effect size, ideally with its 95\,\% confidence interval;
  \item the name of the test and the exact p-value;
  \item in figure captions, what the bars or bands show (SD, SE or 95\,\% CI) and $n$.
\end{enumerate}

For example: ``Class~A scored on average 13.2 marks more than class~C (Welch's $t$-test with Holm's
correction, $p = 0.023$; $n = 18$ per class).'' simplydata's printed answers name each test, and
\code{table.steps()} lists every change you made to the table, in order, for your methods.
).
% Plain words first, then the formal name, so a learner can follow a paper's Methods section.

\unchapter{Describing numbers}{ideas:describe}

Before any test or model, a few numbers describe a column well: where its middle is, and how
spread out it is. simplydata's summary words say each one in a sentence.

\unsection{Average or median?}{ideas:average-median}

The \textbf{average} (formally the \textbf{arithmetic mean}, written $\bar{x}$) adds all the values
and divides by how many there are:
\[
  \bar{x} = \frac{x_1 + x_2 + \dots + x_n}{n}.
\]
The \textbf{median} is the middle value when the values are put in order (with an even number of
values, the average of the two in the middle). Half the values are below it and half above.

Both describe the ``middle'', but they react differently to values far from the rest
(\textbf{outliers}). Five children earn pocket money of 2, 3, 3, 4 and 50 a week: the average is
12.4, though four of the five get 4 or less; the median is 3. One extreme value pulls the average
towards it but hardly moves the median. So:

\begin{itemize}
  \item Use the \textbf{average} when the values are roughly symmetric (a bell shape), and when you
    want to add things up (total = average $\times$ count).
  \item Use the \textbf{median} when the values are lopsided (\textbf{skewed}), such as incomes,
    house prices, waiting times or bacterial counts, or when there are outliers you cannot
    remove.
  \item When the two are far apart, say so: it tells the reader the values are skewed.
\end{itemize}

In simplydata: \code{table.average\_of("marks")}, \code{table.median\_of("marks")},
\code{nums.average()}, \code{nums.median()}, and on charts \code{show="average"} or
\code{show="median"}. Lesson~\ref{lesson:5} draws both.

\unsection{Spread: range, percentiles and standard deviation}{ideas:spread}

Two classes can have the same average mark, one with every child near 65 and the other with
marks from 20 to 100. The \textbf{spread} tells them apart.

\begin{itemize}
  \item The \textbf{range} is the biggest value minus the smallest. It is easy to understand but
    depends on just two values, often the strangest ones.
  \item A \textbf{percentile} is the value that a given share of the values is below:
    \code{nums.percentile(90)} is the value that 90\,\% are below. The median is the 50th
    percentile. The middle half of the values lies between the 25th and 75th percentiles; that
    distance is the \textbf{interquartile range (IQR)}, and it is the box in a box plot
    (\code{chart.boxes}).
  \item The \textbf{standard deviation} (SD, written $s$) is about how far a typical value is from
    the average:
    \[
      s = \sqrt{\frac{\sum_{i=1}^{n} (x_i - \bar{x})^2}{n - 1}}.
    \]
    Its square, $s^2$, is the \textbf{variance}. Dividing by $n-1$ instead of $n$ (Bessel's
    correction) makes the sample variance an unbiased estimate of the variance of everything the
    sample came from; simplydata always divides by $n-1$, as pandas and R do (NumPy divides by $n$
    unless told otherwise).
\end{itemize}

The SD goes with the average and the IQR with the median. For all the marks of the practice table:
``average 64.9 (SD 15.0)'' or ``median 68.5 (IQR 54.3 to 74.8)''.

\unsection{The shape of the values: distributions}{ideas:distribution}

The way values are spread out is their \textbf{distribution}; a histogram
(\code{chart.histogram}) or a density curve (\code{chart.density}) shows it. Many measurements,
such as heights or measuring errors, follow a bell shape called the \textbf{normal} (or
Gaussian) \textbf{distribution}. For a normal distribution about 68\,\% of the values lie within
one SD of the average and about 95\,\% within two SD.

Not everything is bell-shaped. Counts and concentrations are often skewed to the right (many
small values, a few large ones); a log scale (\code{chart.log\_scale\_y()}) often makes them look
more symmetric. Look at the distribution before choosing what to report and which test to use.

\ideachapter{How sure are we?}{ideas:uncertainty}

We almost never measure everything. We measure a \textbf{sample} (18 children, 40 water samples)
and want to learn about the \textbf{population} it came from (all children of that age, the whole
fjord). A number worked out from a sample, such as its average, is an \textbf{estimate}, and
another sample would give a slightly different one. These ideas say how much.

\unsection{Standard deviation or standard error?}{ideas:sd-se}

They sound alike and are often mixed up, but they answer different questions.

\begin{center}\small
\begin{tabularx}{\textwidth}{@{}>{\raggedright\arraybackslash}p{0.19\textwidth}>{\raggedright\arraybackslash}X>{\raggedright\arraybackslash}X@{}}
\toprule
 & \textbf{Standard deviation (SD)} & \textbf{Standard error of the mean (SE, SEM)}\\
\midrule
Describes & the \emph{values}: how spread out they are & the \emph{average}: how precisely the sample pins it down\\
Formula & $s$ (above) & $\mathrm{SE} = s / \sqrt{n}$\\
More data & stays about the same & gets smaller (4 times the data halves it)\\
Use it to & show the spread of individuals & compare averages, build confidence intervals\\
simplydata & \code{error="standard deviation"} & \code{error="standard error"}\\
\bottomrule
\end{tabularx}
\end{center}

Error bars of the SE are always shorter than those of the SD, and can look reassuringly small. That
is why a figure caption must say which one was drawn; simplydata prints a line saying it, ready
for your caption.

\unsection{Confidence intervals}{ideas:ci}

A \textbf{95\,\% confidence interval} (CI) is a range of plausible values for the true number in
the population, worked out from the sample. For an average it is
\[
  \bar{x} \pm t_{0.975,\,n-1} \times \mathrm{SE},
\]
where $t_{0.975,\,n-1}$ is a number from the $t$ distribution: about 2 for most samples, a bit
more for small ones (2.26 for $n = 10$, 2.09 for $n = 20$).

What ``95\,\%'' means, exactly: if you repeated the whole study many times and built an interval
each time, about 95 out of 100 of those intervals would contain the true value. It does
\emph{not} mean that the true value moves around, nor that 95\,\% of the individual values lie
inside (that would be about $\bar{x} \pm 2\,\mathrm{SD}$ for a bell shape).

\begin{itemize}
  \item A \textbf{narrow} interval means the data pin the value down well; a \textbf{wide} one
    means you need more data.
  \item An interval says more than a p-value: it gives the size of the effect \emph{and} how sure
    you are of it. \code{table.fit\_straight\_line} prints the slope with its 95\,\% CI.
  \item On charts: \code{error="confidence interval"} on bars and lines, and
    \code{fit\_line(confidence\_band=True)} for a fitted line.
\end{itemize}

\unsection{Why averages behave well: the central limit theorem}{ideas:clt}

Even when the values themselves are skewed, the \emph{average} of many of them is close to
bell-shaped; this is the \textbf{central limit theorem}. It is why tests on averages (t-tests,
ANOVA) usually work well with 30 or more values per group, even for values that are not
bell-shaped. With small groups of skewed values, rank-based tests are safer
(\code{based\_on="ranks"}).

\ideachapter{Testing: real difference, or chance?}{ideas:testing}

\unsection{The null hypothesis and the p-value}{ideas:p-value}

Two groups almost never have exactly the same average, even when nothing real separates them:
samples differ by chance. A statistical test asks whether a difference is bigger than chance alone
would often produce. It works like a court case, where the accused is innocent until proven
guilty:

\begin{enumerate}
  \item Assume there is \textbf{no real difference}. This assumption is the \textbf{null
    hypothesis} ($H_0$). The \textbf{alternative hypothesis} ($H_1$) is that there is one.
  \item Work out how surprising your data would be if $H_0$ were true. The \textbf{p-value} is the
    probability of a difference at least as big as the one you saw, if there were really no
    difference.
  \item A small p-value means the data are hard to explain by chance alone. By custom, below
    $0.05$ is called \textbf{statistically significant}; that threshold is the
    \textbf{significance level} $\alpha$.
\end{enumerate}

simplydata says the p-value in words: below 0.001 ``very unlikely to be chance'', below 0.05
``unlikely to be chance'', otherwise ``could easily be chance''.

\begin{goodtoknow}
What a p-value is \textbf{not}:
\begin{itemize}
  \item It is \textbf{not} the probability that $H_0$ is true, nor the probability that the result
    is a fluke.
  \item A small p-value does \textbf{not} say the difference is big or important. With thousands
    of values, tiny, meaningless differences become ``significant''. The \emph{effect size} says
    how big it is.
  \item A large p-value does \textbf{not} show that there is no difference. It only says these
    data cannot tell. ``Absence of evidence is not evidence of absence.''
  \item 0.049 and 0.051 are practically the same evidence. Report the exact p-value, not just
    ``$p < 0.05$''.
\end{itemize}
\end{goodtoknow}

\unsection{Two kinds of mistake, and power}{ideas:errors}

\begin{itemize}
  \item A \textbf{false alarm} (\textbf{type I error}, a false positive): the test says there is a
    difference when there is none. With $\alpha = 0.05$ this happens in about 1 test in 20 when
    $H_0$ is true.
  \item A \textbf{missed difference} (\textbf{type II error}, a false negative): there is a real
    difference, but the test does not find it. Small samples miss differences often.
\end{itemize}

The \textbf{power} of a study is the chance of finding a difference of a given size if it is
really there. More data, less noisy measurements and bigger true differences all raise power. A
study with low power can only find large effects; plan the sample size before collecting data.

\unsection{Many tests at once}{ideas:multiple}

Every test at $\alpha = 0.05$ carries a 5\,\% risk of a false alarm, and the risks add up: with 20
tests on pure noise, you expect about one ``significant'' result. Correcting for
\textbf{multiple comparisons} keeps the chance of \emph{any} false alarm at 5\,\%.
\code{table.compare\_each\_pair} uses \textbf{Holm's method}, which is never less powerful than the
simpler Bonferroni correction (multiply each p-value by the number of tests). Decide which
comparisons you will make before you look at the data.

\unsection{Effect sizes: how big is the difference?}{ideas:effect-size}

The p-value answers ``could it be chance?''. The \textbf{effect size} answers ``how big is it?'',
which is usually what matters. Report both, ideally with a confidence interval for the effect.

\begin{center}\small
\begin{tabularx}{\textwidth}{@{}>{\raggedright\arraybackslash}p{0.26\textwidth}>{\raggedright\arraybackslash}X>{\raggedright\arraybackslash}p{0.25\textwidth}@{}}
\toprule
\textbf{Effect size} & \textbf{What it says} & \textbf{Given by}\\
\midrule
Difference in averages & in the units of the data (``7.6 marks higher''): often the clearest & \code{compare\_averages}, \code{compare\_each\_pair}\\
Cohen's $d$, Hedges' $g$ & the difference in standard deviations; $g$ corrects $d$ for small samples. Rough guide: 0.2 small, 0.5 medium, 0.8 large & \code{compare\_averages} (2 groups), \code{compare\_average\_to}, \code{compare\_before\_after} ($d_z$)\\
Omega-squared $\omega^2$ & the share of all the differences that the groups explain (0 to 1) & \code{compare\_averages} (3+ groups)\\
Cramér's $V$ & how strongly two categories are linked, 0 (none) to 1 & \code{compare\_counts}\\
Correlation $r$ & how closely two columns go up and down together, $-1$ to $1$ & \code{correlation\_between}\\
$R^2$ & the share of the differences in $y$ that a line or curve explains & \code{fit\_straight\_line}, \code{fit\_formula}\\
\bottomrule
\end{tabularx}
\end{center}

The ``small, medium, large'' labels (from Jacob Cohen) are only a last resort: a ``small''
difference in survival can be important, and a ``large'' one in a lab assay may not be.

\ideachapter{Links between columns}{ideas:links}

\unsection{Correlation}{ideas:correlation}

A \textbf{correlation} measures how closely two columns of numbers go up and down together, from
$-1$ (one always goes down when the other goes up) through $0$ (no straight-line link) to $1$.

\begin{itemize}
  \item \textbf{Pearson's $r$} (the usual one) measures a \emph{straight-line} link. It is
    sensitive to outliers, and can be near 0 for a strong curved link (a U shape).
  \item \textbf{Spearman's $\rho$} (rho) uses ranks: it measures any link that keeps going one way
    (always up, or always down), and is robust to outliers.
    \code{table.correlation\_between("hours", "marks", method="spearman")}.
  \item $r^2$ is the share of the differences in one column that a straight line through the other
    explains. $r = 0.5$ explains only a quarter.
\end{itemize}

Always draw the dots (\code{chart.dots}) before trusting a correlation: very different patterns
can have the same $r$ (Anscombe's quartet is the famous example).

\unsection{Correlation is not cause}{ideas:cause}

Ice cream sales and drownings go up together, but ice cream does not cause drowning: hot weather
drives both. When two things go together there are several possible reasons:

\begin{enumerate}
  \item A causes B.
  \item B causes A (\textbf{reverse causation}): do people exercise because they are healthy, or
    are they healthy because they exercise?
  \item A third thing causes both (a \textbf{confounder}): hot weather, age, wealth, or the size of
    a hospital.
  \item Chance: test enough pairs of columns and some will line up by luck.
\end{enumerate}

Only an \textbf{experiment} where you change A yourself, ideally by random assignment (a
randomised controlled trial), shows cause directly. With observational data, say ``is associated
with'', not ``causes''. simplydata's link words remind you: ``A link does not show that one thing
causes the other.''

\unsection{Fitting a line or a formula}{ideas:regression}

\textbf{Linear regression} finds the straight line $y = a + b\,x$ that goes as close as possible to
the dots, by making the sum of the squared vertical distances from the dots to the line as small as
possible (\textbf{least squares}). The \textbf{slope} $b$ says how much $y$ changes for each 1 more
of $x$; its p-value tests whether the slope could be 0, and its confidence interval gives a range of
plausible slopes. $R^2$ says how much of the differences in $y$ the line explains.
\code{table.fit\_straight\_line(x="hours", y="marks")} prints all of these;
\code{chart.fit\_line()} draws the line.

When the relation is curved, such as growth that multiplies by the same share each year, write the
formula yourself: \code{table.fit\_formula("population = a * exp(b * (year - 2000))")} finds $a$ and
$b$ by \textbf{nonlinear least squares}. A fitted line or curve describes the data you have;
predicting far outside their range (\textbf{extrapolation}) is risky.

\ideachapter{The shape of a table}{ideas:tables}

\unsection{Tidy data: long and wide tables}{ideas:tidy}

The same marks can be stored in two shapes:

\begin{center}\small
\begin{tabular}{@{}lll@{\hspace{3em}}llll@{}}
\multicolumn{3}{@{}l}{\textbf{Long (tidy)}} & \multicolumn{4}{l@{}}{\textbf{Wide}}\\
\toprule
name & subject & marks & name & Maths & Science & English\\
\midrule
Asha & Maths & 61 & Asha & 61 & 66 & 77\\
Asha & Science & 66 & Ben & 73 & 98 & 71\\
Asha & English & 77 & & & & \\
Ben & Maths & 73 & & & & \\
\multicolumn{3}{@{}l}{\dots} & & & & \\
\bottomrule
\end{tabular}
\end{center}

A table is \textbf{tidy} (Hadley Wickham's term) when each \textbf{variable} is a column, each
\textbf{observation} is a row, and each value has its own cell. The long table is tidy: it has one
column for what was measured (\code{marks}) and columns that say which child and which subject.

\begin{itemize}
  \item \textbf{Long tables} suit charts, groups and tests: \code{chart.bars(x="subject", y="marks")}
    and \code{compare\_averages("marks", between="subject")} need the subject as a column.
  \item \textbf{Wide tables} suit sums across columns (a total per child) and paired tests
    (\code{compare\_before\_after("Maths", "Science")}), and are how people often type data into
    a spreadsheet.
  \item \code{table.turn\_columns\_into\_rows} makes a wide table long (in pandas: \code{melt}; in R:
    \code{pivot\_longer}); \code{table.turn\_rows\_into\_columns} makes a long table wide (\code{pivot}).
\end{itemize}

\unsection{Missing values and imputation}{ideas:missing}

Real data has empty cells: a sample was lost, a question skipped. What to do depends on
\emph{why} they are empty:

\begin{itemize}
  \item \textbf{Missing completely at random}: the reason has nothing to do with the data (a tube
    dropped). Leaving those rows out loses some power but does not bias the result.
  \item \textbf{Missing at random}: the reason depends on other columns you have (older people skip
    a question more often). Leaving rows out can bias the result; taking that other column into
    account helps.
  \item \textbf{Missing not at random}: the reason depends on the missing value itself (very ill
    patients drop out; concentrations below the detection limit). No simple fix removes the bias;
    say so in your report.
\end{itemize}

\textbf{Removing} rows with empty cells (\code{table.remove\_rows\_with\_empty\_cells()}) is
simple and honest when few rows are affected. \textbf{Filling} them in is called
\textbf{imputation}: \code{table.fill\_empty\_cells("marks", with\_the="average")} (also
\code{"median"}, \code{"value above"}, \code{"value below"}, or \code{with\_value=0}). Filling with
the average keeps the average the same, but makes the values look \emph{less spread out} than they
really are, so standard deviations, standard errors and p-values become too small. Filling with the
value above or below suits a series in time order. Whatever you choose, count the empty cells first
(\code{table.count\_empty\_cells()}) and say in your methods how you handled them.

\ideachapter{Which test when?}{ideas:which-test}

First draw the data: a chart shows outliers, skew and odd values that a single p-value hides. Then
answer four questions:

\begin{enumerate}
  \item \textbf{What is the question?} A difference between groups, a change within the same
    rows, a link between two columns, or the shape of a relation?
  \item \textbf{What kind of data?} Numbers (marks, heights, concentrations) or categories (pass or
    fail, species, site)?
  \item \textbf{How many groups, and are they the same rows?} Two groups or more; independent
    groups (different children) or paired measurements (the same children twice)?
  \item \textbf{Are the groups small and the values far from bell-shaped?} Then use the rank-based
    version: \code{based\_on="ranks"}. simplydata checks each group for a bell shape
    (Shapiro--Wilk test) and adds a note when it matters.
\end{enumerate}

\begin{center}\footnotesize
\renewcommand{\arraystretch}{1.25}
\begin{tabularx}{\textwidth}{@{}>{\raggedright\arraybackslash}p{0.25\textwidth}>{\raggedright\arraybackslash}p{0.27\textwidth}>{\raggedright\arraybackslash}X>{\raggedright\arraybackslash}X@{}}
\toprule
\textbf{Question} & \textbf{simplydata word} & \textbf{Test (the usual one)} & \textbf{With \code{based\_on="ranks"}}\\
\midrule
Is one average different from a number you choose? & \code{compare\_average\_to("marks", 60)} & one-sample $t$-test & Wilcoxon signed-rank test\\
Do \textbf{two} groups differ in their average? & \code{compare\_averages("marks", between="group")} & Welch's $t$-test & Mann--Whitney $U$ test\\
Do \textbf{three or more} groups differ? & \code{compare\_averages("marks", between="class")} & Welch's one-way ANOVA & Kruskal--Wallis test\\
\emph{Which} groups differ? & \code{compare\_each\_pair("marks", between="class")} & pairwise Welch $t$-tests, Holm's correction & pairwise Mann--Whitney tests, Holm's correction\\
Did the values change, in the \textbf{same} rows measured twice? & \code{compare\_before\_after("before", "after")} & paired $t$-test & Wilcoxon signed-rank test\\
Do \textbf{counts} of categories differ between groups? & \code{compare\_counts("result", between="class")} & chi-square test of independence & --- (for a $2\times2$ table with small counts: \code{test="fisher"}, Fisher's exact test)\\
Do two columns of numbers go up and down together? & \code{correlation\_between("hours", "marks")} & Pearson's $r$ & \code{method="spearman"}: Spearman's $\rho$\\
How much does $y$ change with $x$? & \code{fit\_straight\_line(x="hours", y="marks")} & linear regression (least squares) & ---\\
Does the data follow a formula (growth, decay, saturation)? & \code{fit\_formula("y = a * exp(b * x)")} & nonlinear least squares & ---\\
\bottomrule
\end{tabularx}
\end{center}

\begin{goodtoknow}
\begin{itemize}
  \item \textbf{Welch, not Student.} For averages simplydata uses Welch's versions, which do not
    assume that the groups have the same spread. They are R's defaults and lose almost nothing when
    the spreads are equal. An expert can still choose: \code{test="student"} or \code{test="anova"}.
  \item \textbf{Paired or not?} If each row of one measurement belongs to a row of the other (the same
    child, the same site), the test must be paired: \code{compare\_before\_after}. Treating paired
    data as independent groups throws away information.
  \item \textbf{Counts, not percentages.} \code{compare\_counts} needs the rows themselves (one row per
    observation), not percentages worked out by hand.
  \item \textbf{One test per question.} Decide the question and the test before looking at the
    results; trying several tests and keeping the best p-value makes false alarms much more
    likely.
  \item Each test gives back a one-row table with every number in it:
    \code{result = table.compare\_averages(...)}, then \code{result.save("test.csv")}.
\end{itemize}
\end{goodtoknow}

\unsection{Writing it up}{ideas:report}

A good results sentence lets a reader check your work. Give:
\begin{enumerate}
  \item the summary for each group (average and SD, or median and IQR) and the number of values $n$;
  \item the effect size, ideally with its 95\,\% confidence interval;
  \item the name of the test and the exact p-value;
  \item in figure captions, what the bars or bands show (SD, SE or 95\,\% CI) and $n$.
\end{enumerate}

For example: ``Class~A scored on average 13.2 marks more than class~C (Welch's $t$-test with Holm's
correction, $p = 0.023$; $n = 18$ per class).'' simplydata's printed answers name each test, and
\code{table.steps()} lists every change you made to the table, in order, for your methods.
